\documentclass[11pt,a4paper]{article}

\usepackage[utf8]{inputenc}
\usepackage[T1]{fontenc}
\usepackage[margin=2.5cm]{geometry}
\usepackage{booktabs}
\usepackage{tabularx}
\usepackage{enumitem}
\usepackage{graphicx}
\usepackage{hyperref}

\title{Roomie: User Feedback Report}
\author{Mariia Denysova \and Anna Liuklian \and Sofiia Rubakha}
\date{}

\begin{document}

\maketitle

\section{Introduction}
Roomie is a simple web app that helps students living together divide cleaning,
shared shopping and household chores fairly. This report describes who we showed
the app to, what they said, and how the current app addresses their feedback.

The first version of Roomie lets users create a private roommate group and invite
members by link, add weekly chores with assignees and due dates, configure chore
rotation, view weekly chore progress, keep a shared shopping list, and receive
in-app notifications about events such as chore assignments. We aimed to show it to five roommate groups from a university
dormitory.

App link: \url{https://your-app-link-here}

\section{Who We Showed It To}
We showed Roomie to \textbf{17} people from \textbf{5} roommate groups.

\begin{table}[h]
\centering
\begin{tabularx}{\textwidth}{l l l X}
\toprule
\textbf{Group} & \textbf{Members} & \textbf{Living situation} & \textbf{Notes} \\
\midrule
Group 1 & 3 & dorm & Block-style dorm unit; shares a bathroom and small kitchenette \\
Group 2 & 4 & dorm & First-year roommates in one room; previously relied on verbal agreements \\
Group 3 & 3 & apartment & Renting a 2-bedroom flat near campus; struggled to track shared supplies \\
Group 4 & 4 & dorm & Upper-year students; frequent disputes over taking out trash and sweeping \\
Group 5 & 3 & apartment & Busy class and work schedules; needed clear deadlines for weekend cleaning \\
\bottomrule
\end{tabularx}
\caption{Groups who tried the app}
\end{table}

\section{How Many People Tried It}
\begin{table}[h]
\centering
\begin{tabularx}{\textwidth}{X c c}
\toprule
\textbf{Measure} & \textbf{Our target} & \textbf{Actual} \\
\midrule
People who saw the app & --- & 17 \\
People who actually used it & --- & 14 \\
Groups that created a schedule & at least 4 & 5 \\
Groups that marked chores as completed & at least 3 & 4 \\
Groups that kept using Roomie & at least 2 & 3 \\
\bottomrule
\end{tabularx}
\caption{Usage compared with the targets from our proposal}
\end{table}

\section{What They Said}

\subsection{What users liked}
\begin{itemize}[noitemsep]
  \item Weekly chore rotation and the personal chore filter made personal responsibilities immediately visible, reducing awkward reminders in group chats.
  \item The shared shopping list prevented roommates from accidentally buying duplicate household staples like dish soap and paper towels.
  \item Joining a household via invite link took less than a minute and required no complicated setup.
\end{itemize}

\subsection{What users found confusing or missing}
\begin{itemize}[noitemsep]
  \item There was no quick way to swap a chore when a roommate was sick or away for the weekend.
  \item Roommates wanted reminders before larger tasks like bathroom or kitchen cleaning were due.
\end{itemize}

\subsection{Selected quotes}
\begin{quote}
``It saved us from having that passive-aggressive `whose turn is it to buy dish soap?' message in our group chat every week.'' --- Second-year student, dorm resident
\end{quote}
\begin{quote}
``The automatic weekly rotation and `My Tasks' view make it super easy to see what I need to do, though I wished there was a quick button to swap shifts when I go home for the weekend.'' --- Third-year student, shared apartment resident
\end{quote}

\section{How the Current App Addresses Feedback}

\begin{table}[h]
\centering
\begin{tabularx}{\textwidth}{X X X}
\toprule
\textbf{Feedback} & \textbf{Current app} & \textbf{Possible improvement} \\
\midrule
Roommates could not easily handle exceptions when someone went home for the weekend. & Roommates can manually change the assignee or due date of an incomplete chore. & Add a dedicated action to swap upcoming chores. \\
Roommates wanted advance notice for larger weekly chores. & Chores display due dates, and the in-app bell shows events such as new assignments. & Add scheduled reminders before chores are due. \\
\bottomrule
\end{tabularx}
\caption{Feedback and possible improvements}
\end{table}

\section{What We Did Not Change}
\begin{itemize}[noitemsep]
  \item \textbf{In-app group messaging} --- Roommates already coordinate daily in Telegram or WhatsApp, so building a separate chat inside Roomie would add unnecessary clutter.
  \item \textbf{Native iOS and Android apps} --- Developing separate mobile apps was outside our timeline, so we focused on making the web interface responsive on mobile browsers.
\end{itemize}

\section{Bonus: Willingness to Pay}
To check whether users value Roomie enough to pay for it, we provided a link to a
charity fund (a ``jar'') for a volunteer cause. Users were free to contribute, and
people contributed money to it.

\textbf{Charity fund:} [name of the fund / cause] \\
\textbf{Link:} \url{https://your-jar-link-here}

\begin{table}[h]
\centering
\begin{tabularx}{0.8\textwidth}{X c}
\toprule
\textbf{Measure} & \textbf{Value} \\
\midrule
People who were shown the link & [number] \\
People who contributed & [number] \\
Total amount collected & [amount] \\
Average contribution & [amount] \\
\bottomrule
\end{tabularx}
\caption{Contributions to the charity fund}
\end{table}

\subsection{Evidence}
\begin{figure}[h]
\centering
% Replace the box below with your screenshot of the jar, for example:
% \includegraphics[width=0.8\textwidth]{jar_screenshot.png}
\fbox{\parbox{0.8\textwidth}{\centering\vspace{3em}[Screenshot of the jar showing contributions]\vspace{3em}}}
\caption{Charity fund page showing contributions from users}
\end{figure}

\subsection{Comments from contributors}
\begin{itemize}[noitemsep]
  \item {[Why they decided to contribute]}
  \item {[Another comment]}
\end{itemize}

\section{Conclusion}
We showed Roomie to roommate groups to see whether a shared schedule for chores and
shopping is useful in real shared homes. Participants highlighted the value of seeing their own assignments and sharing
a shopping list. They also asked for quicker chore handoffs and advance notice
before chores are due. The current app supports manual reassignment and in-app
assignment notifications. Next, we would explore a dedicated swap action and
scheduled reminders, with Telegram or push notifications as possible delivery
options.

\appendix
\section{Questions We Asked Users}
\begin{enumerate}[noitemsep]
  \item How do you and your roommates currently divide cleaning and chores?
  \item Was it clear how to create a household and invite your roommates?
  \item Was it easy to add a chore, assign it and mark it as done?
  \item Was anything confusing or missing in the app?
  \item Which feature would you use the most (chores, weekly rota, shopping list)?
  \item Would you use Roomie with your roommates? Why or why not?
  \item Would you contribute to the charity fund as a way to support the app? Why or why not?
\end{enumerate}

\end{document}
